\documentclass[../../../include/open-logic-section]{subfiles}

\begin{document}

\olfileid{sth}{spine}{valpha}

\olsection{దశలను $V_\alpha$లుగా నిర్వచించడం}

\olref[sth][ordinals][]{chap}లో
సుక్రమాలనూ (వాన్
న్యూమన్) క్రమసంఖ్యలనూ
నిర్వచించాం. ఈ
అధ్యాయంలో వాటితో
సమితుల స్థాయిక్రమాన్ని
\emph{స్వయంగా} వివరించబోతున్నాం.
దీని కోసం
\olref[sth][ordinals][opps]{sec}లో
ఉత్తరవర్తి, సీమా
క్రమసంఖ్యల భావాలను
నిర్వచించామని
గుర్తుచేసుకోండి. ఈ
భావాలను కింది
నిర్వచనంలో వాడతాం:

\begin{defn}\ollabel{defValphas}
	\begin{align*}
	V_\emptyset &\defis \emptyset\\
	V_{\ordsucc{\alpha}} &\defis \Pow{V_\alpha} & & 
	\alpha\text{ ఏ క్రమసంఖ్య అయినా}\\
	V_{\alpha} &\defis \bigcup_{\gamma < \alpha} V_\gamma & & 
	\alpha\text{ సీమా క్రమసంఖ్య అయినప్పుడు}
\end{align*}
\end{defn}
\noindent
ఇది క్రమసంఖ్యలపై
\emph{అతిపరిమిత పునరావృత్తి}
ద్వారా ఇచ్చే నిర్వచనం.
ఈ విషయంలో దీనిని
సహజ సంఖ్యలపై
ప్రమేయాల పునరావృత్త
నిర్వచనాలతో పోల్చాలి.\footnote{\olref[sfr][infinite][dedekind]{sec}లోని
సంకలనం, గుణకారం,
ఘాతాంక నిర్వచనాలతో
పోల్చండి.} సహజ సంఖ్యల
విషయంలో మాదిరిగానే
ఇక్కడా ప్రాథమిక సందర్భం,
ఉత్తరవర్తి సందర్భాలను
నిర్వచిస్తాం; కానీ
క్రమసంఖ్యల విషయంలో
\emph{సీమా} సందర్భాల
ప్రవర్తనను కూడా
వివరించాలి.

$V_\alpha$ల ఈ
నిర్వచనం ఒక ముఖ్యమైన
మైలురాయి అవుతుంది.
మన సమితి స్థాయిక్రమంలో
\emph{దశలవారీగా}
సమితులు ఏర్పడతాయని
అనౌపచారికంగా
సమర్థించాం. $V_\alpha$లు
నిజానికి ఆ దశలే.
అయితే ఇది దశల
\emph{అంతర్గత} వివరణ
అనే విషయం ముఖ్యం.
సిద్ధాంతానికి ఒక సాధ్యమైన
\emph{నమూనా}ను సూచించడానికి
బదులు, మన సమితి
సిద్ధాంతం \emph{లోపలే}
దశలను నిర్వచిస్తాం.

\end{document}
